\section{El significado estricto de Open Source: de «utilizable» a «reproducible»}

\subsection{La naturaleza técnica de la disputa sobre la definición}
En la segunda mitad, el ponente discutió el problema de los límites de la definición de apertura: si solo se cuenta con los pesos y faltan los datos de entrenamiento y el proceso completo, ¿basta eso para llamarlo open source? No se trata de una disputa semántica, sino de algo que afecta la reproducibilidad de la investigación científica y la delimitación de responsabilidades.

\begin{importantbox}{Open Source for Models necesita un nuevo marco de definición}
El open source tradicional del software gira principalmente en torno a que «el código fuente sea visible y modificable».  
En los sistemas de modelos, los datos, el proceso de entrenamiento, el protocolo de evaluación y la estrategia de publicación son igualmente componentes de la «fuente».
\end{importantbox}

\subsection{La escalera de la reproducibilidad}
\begin{table}[H]
\centering
\begin{tabular}{p{3.0cm}p{3.8cm}p{4.8cm}}
\toprule
\textbf{Nivel de apertura} & \textbf{Reproducción posible} & \textbf{Lo que aún falta} \\
\midrule
API disponible & Reproducción del comportamiento, evaluación a nivel de interfaz & Atribución del mecanismo, rastreo del entrenamiento \\
Weights disponibles & Reproducción de la inferencia, reproducción del fine-tuning, experimentos parciales de interpretación & La cadena causal de datos y entrenamiento \\
Full source & Reproducción de extremo a extremo del entrenamiento y la evaluación, modificación a nivel de sistema & El principal cuello de botella pasa a ser el cómputo y los recursos de ingeniería \\
\bottomrule
\end{tabular}
\caption{La escalera de capacidades de «poder invocarse» a «poder reproducirse»}
\end{table}

\begin{knowledgebox}{Por qué «reproducible» es más importante que «descargable»}
La comunidad científica necesita una acumulación de conocimiento verificable. Solo cuando las condiciones experimentales clave son reproducibles, comparables y refutables, el campo puede formar un consenso sólido, en lugar de depender de las afirmaciones no verificables de organizaciones individuales.
\end{knowledgebox}

\begin{warningbox}{Tratar «open source» como una etiqueta de mercadotecnia termina por dañar la confianza de la comunidad}
Si el término se usa mal de manera sostenida, los investigadores cometerán errores sistemáticos en sus expectativas de colaboración:  
creerán que algo es reproducible cuando en realidad no lo es; creerán que algo admite auditoría cuando en realidad no la admite.
\end{warningbox}

\subsection{Resumen de la sección}
El núcleo de esta sección es que «la definición es infraestructura». Para los sistemas de modelos, la definición de apertura determina directamente la reproducibilidad de la investigación, la eficiencia de la colaboración y la operatividad de la gobernanza.


